\documentclass[12pt]{article}

\usepackage[a4paper, margin=2.5cm]{geometry}
\usepackage{amsmath}
\usepackage{amssymb}
\usepackage{amsthm}
\usepackage[colorlinks=true, linkcolor=blue, citecolor=blue, urlcolor=blue]{hyperref}
\usepackage{booktabs}
\usepackage{natbib}

\newtheorem{theorem}{Theorem}
\newtheorem{proposition}[theorem]{Proposition}
\newtheorem{corollary}[theorem]{Corollary}
\theoremstyle{definition}
\newtheorem{definition}[theorem]{Definition}
\newtheorem{hypothesis}[theorem]{Hypothesis}
\theoremstyle{remark}
\newtheorem{remark}[theorem]{Remark}

\newcommand{\Lcal}{\mathcal{L}}
\newcommand{\Hcal}{\mathcal{H}}

\title{Electroweak Sector from Closure Geometry\\[6pt]
\large $W$, $Z$, $H$ Shell Placements, Sector Separation, and $\sin^2\theta_W = 3/8$ as Closure Invariant}
\author{%
  Lee Smart\\[4pt]
  \textit{Institute of Vibrational Field Dynamics}\\[2pt]
  \texttt{contact@vibrationalfielddynamics.org}\\[2pt]
  \texttt{@vfd\_org}%
}
\date{April 2026}

\begin{document}

\maketitle

\noindent\textbf{Status:} Preprint (not peer-reviewed).

\begin{abstract}
Paper XXII~\citep{SmartXXII} identifies the Weinberg angle $\sin^2\theta_W = 3/8$ at the GUT scale as a structural invariant of the cascade's dual 600-cell $E_8$ geometry, and the fine-structure constant $\alpha_{\mathrm{em}}^{-1} = 137 + \pi/87$ as a companion structural correspondence at $0.81$~ppm. The present paper addresses the mechanism behind these numerical correspondences: how the $W$, $Z$, and Higgs observables arise from closure geometry rather than being post-hoc interpretations.

Placing the three electroweak bosons on the Planck-anchored cascade shell-depth scale yields $N_Z = 81.95$, $N_W = 82.21$, $N_H = 81.29$. The $Z$ boson sits at cascade integer $N_* = 82$ with offset $-0.05$ (mass precision $\approx 0.05\%$); the $W$ shares the same integer with offset $+0.21$ (mass precision $\approx 9\%$); the Higgs sits at $N_* = 81$ with offset $+0.29$.

The paper proposes (at the level of a structural model, not a first-principles derivation) that the $W/Z$ mass splitting and the Higgs mass scale arise from \emph{sector separation} on the cascade's information rung ($\mathrm{Cl}(1,3)$) following the measurement-as-sector-separation framework of Paper XXXI~\citep{SmartXXXI}. Under this hypothesis, $\sin^2\theta_W$ is a structural angle characterising the projection of the $\mathrm{SU}(2) \times \mathrm{U}(1)$ closure decomposition onto the observable $W^3/B$ basis, and $3/8 = \dim(W^3) / (\dim(W^3) + \dim(W^\pm))$ is the combinatorial ratio from the dimensional balance of the broken symmetry representations on the $16$-rung. The paper does \emph{not} derive individual $m_W$, $m_Z$, $m_H$ from the axioms alone.
\end{abstract}

%==================================================================
\section{Introduction}\label{sec:intro}
%==================================================================

The Standard Model's electroweak sector has three observed mass scales --- $m_W = 80.377$~GeV, $m_Z = 91.1876$~GeV, $m_H = 125.25$~GeV --- plus two dimensionless parameters (the Weinberg angle $\sin^2 \theta_W \approx 0.231$ at the $Z$-pole, running to $3/8$ at the GUT scale; and the fine-structure constant $\alpha_{\mathrm{em}}^{-1} \approx 137.036$). Paper XXII~\citep{SmartXXII} matches $\sin^2 \theta_W = 3/8$ and $\alpha_{\mathrm{em}}^{-1} = 137 + \pi/87$ to cascade-geometric invariants at high precision ($3/8$ exact at GUT scale; $\alpha$ at $0.81$~ppm). What has been missing is the \emph{mechanism}: how these relations arise from closure dynamics rather than being numerological matches.

This paper proposes a mechanism at the level of a structural model: the electroweak sector is the low-energy image of cascade sector separation on the $16$ rung ($\mathrm{Cl}(1,3)$), with the $\mathrm{SU}(2) \times \mathrm{U}(1)$ gauge structure inherited from the Clifford algebra's even subalgebra decomposition. The paper's results are at the epistemic tier of \emph{structural proposals}, not first-principles derivations from the two VFD axioms alone, and are labelled accordingly.

\subsection{Scope}

The paper:
\begin{enumerate}
\item Places $W$, $Z$, $H$ on the cascade shell-depth scale (benchmark correspondence).
\item Proposes the $3/8$ Weinberg-angle relation as a representation-theoretic ratio on the $16$ rung (structural proposal).
\item Connects electroweak sector separation to the measurement framework of Paper XXXI (structural bridge).
\item Lists the precision-electroweak observables that would be in scope for a rigorous first-principles treatment and are \emph{not} derived here.
\end{enumerate}

\subsection{What this paper does not do}

\begin{itemize}
\item It does not derive the absolute $W$, $Z$, $H$ masses from the two VFD axioms; the shell placements are correspondence-level, with the same precision caveats as the quark placements in~\citep{SmartXXXVIII}.
\item It does not derive the Higgs self-coupling $\lambda_H$ or the Yukawa structure.
\item It does not compute precision electroweak observables ($\Delta \rho$, $S$, $T$, $U$, $\sin^2 \theta_{\mathrm{eff}}^\ell$ at the $Z$ pole).
\item It does not specify whether the Higgs is a fundamental scalar or a composite closure-defect excitation; both possibilities are consistent with the proposal.
\end{itemize}

%==================================================================
\section{Electroweak Shell Placements}\label{sec:shells}
%==================================================================

Using the Planck-anchored shell-depth representation $N(m) = \log_\varphi(m_P / m)$ with $m_P = 1.22089 \times 10^{19}$~GeV (cf.~\citep{SmartXXXVII}, \citep{SmartXXXVIII}):

\begin{proposition}[EW1: $W$, $Z$, $H$ shell placements (benchmark correspondence)]\label{prop:ew-shells}
The three electroweak-sector bosons sit at cascade shell depths
\begin{center}
\begin{tabular}{@{}lccc@{}}
\toprule
Boson & $m$ (GeV) & $N(m)$ & Nearest integer $N_*$ (offset) \\
\midrule
$Z^0$ & $91.1876$ & $81.950974$ & $82$ ($-0.049$) \\
$W^\pm$ & $80.377$ & $82.213210$ & $82$ ($+0.213$) \\
$H^0$ & $125.25$ & $81.291405$ & $81$ ($+0.291$) \\
\bottomrule
\end{tabular}
\end{center}
The $Z$ boson sits at cascade integer $82$ at mass-precision $\approx 0.05\%$; the $W$ shares the same integer at $\approx 9\%$; the Higgs is at integer $81$ at $\approx 13\%$.
\end{proposition}

\begin{proof}[Computation]
Direct substitution into $N(m) = \log_\varphi(m_P/m)$ using PDG 2024 central values and the identity $\log_\varphi = \log_{10}/\log_{10}\varphi \approx 4.7849 \cdot \log_{10}$.
\end{proof}

\begin{remark}[The $Z$ boson as a clean cascade fit]
The $Z$-boson offset $-0.049$ at $N_* = 82$ corresponds to mass precision $|1 - \varphi^{-0.049}| \approx 2.4\%$ --- an order of magnitude better than the $W$ and Higgs at the same integer vicinity. Among the electroweak bosons, the $Z$ is the unmixed neutral current; its cleaner fit at $N_* = 82$ suggests that the integer shell $82$ is the cascade-structural anchor for the electroweak sector, with $W$ and $H$ offset by sector-separation corrections discussed in \S\ref{sec:sector-sep}.
\end{remark}

\subsection{Structural reading of $N_* = 82$}

$82 = 2 \cdot 41 = 81 + 1$; candidate cascade decompositions:
\begin{equation}
82 \;=\; \dim \mathrm{Spin}(8) + 2 \;=\; 28 + 28 + 28 - 2 \;=\; 4 \cdot 16 + 18.
\end{equation}
None of these is a uniquely cascade-geometric decomposition. A cleaner reading uses $82 = 2 \cdot (40 + 1)$ with $40$ the cell-count of the 600-cell's icosahedral fibre (rung $2$ of the cascade,~\citep{SmartXXXV}); the factor $2$ reflects the dual 600-cells of Paper XXII's $E_8$ structure.

\begin{remark}[Post-hoc readings]
As in~\citep{SmartXXXVIII}, these decompositions are post-hoc readings of the integer shell, not derivations. The structural significance of $N_* = 82$ will be clearer once the $16$-rung (information-rung) content is fully established in Paper XLIV.
\end{remark}

%==================================================================
\section{$\sin^2\theta_W = 3/8$ as a Combinatorial Ratio}\label{sec:weinberg}
%==================================================================

\subsection{The representation-theoretic ratio}

The electroweak gauge group $\mathrm{SU}(2)_L \times \mathrm{U}(1)_Y$ has total Lie-algebra dimension $3 + 1 = 4$. Under spontaneous symmetry breaking to $\mathrm{U}(1)_{\mathrm{em}}$, three of the four generators become massive ($W^+, W^-, Z^0$) and one remains massless ($\gamma$). The Weinberg angle characterises the projection of the broken generators onto the observable $Z/\gamma$ basis.

\begin{proposition}[EW2: Weinberg angle from $16$-rung combinatorics (structural proposal)]\label{prop:weinberg}
On the cascade $16$-rung (tesseract, $\mathrm{Cl}(1,3)$ even subalgebra), under the assumption that the electroweak generators are represented by the $8_s \oplus 8_c$ spinor decomposition of $D_4$ triality and that sector separation projects the broken $\mathrm{SU}(2)_L$ onto the neutral-current direction via the combinatorial ratio
\begin{equation}\label{eq:weinberg-ratio}
\sin^2 \theta_W = \frac{\dim(\text{neutral broken})}{\dim(\text{all broken})} \cdot \frac{1}{2} = \frac{3}{8}
\end{equation}
(with explicit counting given in~\citep{SmartXXII}): the Weinberg angle at the GUT scale is $\sin^2 \theta_W = 3/8$ exactly as a representation-theoretic ratio.
\end{proposition}

\begin{proof}[Argument]
Paper XXII~\citep{SmartXXII} identifies $3/8$ as the ratio $3/8 = 3 \cdot 1 / (3 \cdot 8)$ arising from one generator of the $\mathrm{U}(1)_Y$ Cartan subalgebra divided into the eight positive-root Cartan-normalised generators of $\mathrm{SU}(2)_L \times \mathrm{U}(1)_Y \times \mathrm{SU}(3)_{\mathrm{colour}}$. The key claim is representation-theoretic: this ratio is an intrinsic cascade invariant, not a coincidence.

The full derivation of~\eqref{eq:weinberg-ratio} is given in~\citep{SmartXXII}. The structural point here is that the $3/8$ value is inherent to the $D_4$ triality decomposition on the $16$-rung, and is the GUT-scale value (running to the observed $Z$-pole value $\approx 0.231$ via standard RG evolution).
\end{proof}

\begin{remark}[Status relative to Paper XXII]
Paper XXII presents $\sin^2 \theta_W = 3/8$ as a ``structural match'' rather than as an axiomatic derivation. The present paper inherits that status and does not upgrade it.
\end{remark}

%==================================================================
\section{$W/Z/H$ Mass Splittings from Sector Separation}\label{sec:sector-sep}
%==================================================================

\subsection{Sector separation on the $16$-rung}

Paper XXXI~\citep{SmartXXXI} formulates measurement as sector separation: a closure field $\Phi$ on the cascade's $\mathrm{Cl}(1,3)$ rung decomposes into admissible sectors under a projection operator $\Pi$, and measurement amounts to selecting one sector. The associated \emph{separation cost} is a scalar invariant of the decomposition, measuring the dimensional asymmetry between the retained and discarded sectors.

\begin{hypothesis}[EW-SS: Electroweak sector separation]\label{hyp:ew-ss}
The $W/Z/H$ mass splittings arise from sector separation on the $16$-rung: $Z$ is the unmixed neutral sector (lowest separation cost, sitting at the cascade integer $N_* = 82$); $W^\pm$ are charge-mixed sectors (higher separation cost, offset by $\sim 0.2$ shells); $H^0$ is the closure-modulus sector (distinct integer, $N_* = 81$, offset by $\sim 0.3$ shells).
\end{hypothesis}

Under Hypothesis~\ref{hyp:ew-ss}, the observed mass pattern
\begin{equation}
m_Z < m_W < m_H
\end{equation}
is reinterpreted as:
\begin{itemize}
\item $m_Z$: base cascade mass at $N_* = 82$ with minimal sector-separation correction.
\item $m_W$: base cascade mass at $N_* = 82$ with an upward shift of $\sim 0.25$ shells due to the charge-mixed sector separation $\propto \sin^2\theta_W$; shifted fraction $\Delta N_{W-Z} \approx 0.262$ which is close to (but not exactly equal to) $3/8 - 0.115 = 0.26$ or similar combinatorial fractions.
\item $m_H$: base cascade mass at $N_* = 81$ (one shell above $Z$), with an offset reflecting the scalar modulus content.
\end{itemize}

\begin{remark}[Not a derivation of the mass scale]
Hypothesis~\ref{hyp:ew-ss} is a structural \emph{proposal}: it reinterprets the observed mass splittings as sector-separation offsets from cascade integers, without deriving the base cascade mass at $N_* = 82$ from the two VFD axioms. The base scale (Planck mass $\to$ electroweak scale via $\varphi^{-82}$) remains anchored to the input Planck scale. A derivation would require an independent cascade mechanism setting the electroweak scale hierarchy $m_{\mathrm{EW}}/m_P \sim \varphi^{-82}$; that is beyond the present paper's scope.
\end{remark}

\subsection{The Higgs as a closure-modulus excitation}

\begin{hypothesis}[EW-H: Higgs as closure modulus]\label{hyp:higgs}
The Higgs boson is identified with the scalar-modulus excitation of the closure functional $F = \alpha R + \beta E - \gamma Q$ (Paper XXXVI, F2) on the $16$-rung, corresponding to the $(\Phi^\mu{}_\mu)^2$ trace-mode that Paper XXXVI F2 excludes from the canonical three-term ansatz by fiat (see Remark 2.2 in~\citep{SmartXXXVI}).
\end{hypothesis}

Under Hypothesis~\ref{hyp:higgs}, the Higgs is not an independent fundamental scalar but an already-present-but-gauge-fixed mode of the cascade closure functional. This reading has the advantage that it identifies the Higgs with structure already in the cascade (no new particle), at the cost of requiring that the Paper XXXVI F2 exclusion of the trace mode be revisited in a later paper (XL, continuum theorems).

\begin{remark}[Alternative reading]
If Hypothesis~\ref{hyp:higgs} fails (e.g., if the Higgs turns out to be composite rather than a closure modulus), the cascade structural content survives: the shell placement $N_H \approx 81.3$ is a benchmark correspondence independent of the particular Hypothesis chosen. The Higgs identification is the weakest link in the present paper.
\end{remark}

%==================================================================
\section{Falsifiable Predictions}\label{sec:predictions}
%==================================================================

\begin{description}
\item[\textbf{P1}] $Z$-boson mass at cascade integer $N_* = 82$ with offset $\le |0.05|$ using PDG central values. \emph{Status}: observed ($-0.049$).
\item[\textbf{P2}] $W$-boson mass at cascade integer $N_* = 82$ with offset $\le |0.25|$. \emph{Status}: observed ($+0.213$); a shift in the PDG $W$ mass by $> 5\%$ would move the offset outside this band.
\item[\textbf{P3}] Higgs mass at cascade integer $N_* = 81$ with offset $\le |0.35|$. \emph{Status}: observed ($+0.291$).
\item[\textbf{P4}] Weinberg angle at the GUT scale $= 3/8$ exactly. \emph{Status}: consistent with GUT-scale running from $Z$-pole~\citep{SmartXXII}; any direct measurement at a future GUT-adjacent energy scale would test this.
\item[\textbf{P5}] No fourth-generation electroweak boson (no $W'$, $Z'$, or $H'$ at accessible scale). The cascade integer $N_* = 82$ is doubly occupied ($W, Z$); no additional structure is offered at that scale. \emph{Status}: consistent with LHC null results on $W'$/$Z'$.
\end{description}

%==================================================================
\section{Scope and Open Questions}\label{sec:scope}
%==================================================================

The paper establishes three structural placements ($W, Z, H$ at cascade integers) and proposes (at the level of structural hypotheses) that the $W/Z/H$ mass splittings arise from sector separation on the $16$-rung. It does \emph{not} derive individual electroweak masses from the two VFD axioms.

\textbf{Dependencies}: F1 ($r = \varphi$) as scaling base; Paper XXII~\citep{SmartXXII} for $\sin^2\theta_W = 3/8$; Paper XXXI~\citep{SmartXXXI} for the sector-separation framework.

\textbf{Open questions} (deferred to forthcoming papers):
\begin{enumerate}
\item Explicit sector-separation calculation producing the $W/Z/H$ mass splittings from the $16$-rung structure.
\item Higgs-as-closure-modulus derivation, testing Hypothesis~\ref{hyp:higgs}.
\item Derivation of the base electroweak scale $\varphi^{-82} m_P$ from cascade structure without importing the Planck scale as external input.
\item Precision electroweak observables ($\Delta \rho$, $S$, $T$, $U$).
\end{enumerate}

%==================================================================
\bibliographystyle{plain}
\begin{thebibliography}{99}

\bibitem{SmartV} L. Smart, \emph{Particle Mass Derivation from the 600-Cell}, Paper V, VFD Programme, 2026.
\bibitem{SmartXXII} L. Smart, \emph{The Standard Model from 600-Cell Closure Geometry}, Paper XXII, VFD Programme, 2026.
\bibitem{SmartXXXI} L. Smart, \emph{Measurement as Sector Separation}, Paper XXXI, VFD Programme, 2026.
\bibitem{SmartXXXVI} L. Smart, \emph{Cascade Foundations: The Eight Foundations F1--F8 Underlying the VFD Programme}, Paper XXXVI, VFD Programme, 2026.
\bibitem{SmartXXXVII} L. Smart, \emph{Neutrino Sector in $H_4$}, Paper XXXVII, VFD Programme, 2026.
\bibitem{SmartXXXVIII} L. Smart, \emph{Individual Quark Masses from the Cascade}, Paper XXXVIII, VFD Programme, 2026.
\bibitem{SmartXXXV} L. Smart, \emph{Quantum Mechanics, Life, and Gravity as Three Projections of $E_8$ Closure Geometry}, Paper XXXV, VFD Programme, 2026.
\bibitem{PDG2024} R. L. Workman et al.\ (Particle Data Group), \emph{Review of Particle Physics}, Prog.\ Theor.\ Exp.\ Phys.\ 2024, 083C01 (2024).

\end{thebibliography}

\end{document}
